\section{Identifying Policy Positions in Political Manifestos}
\label{sec:v8-manifesto-mechanism}

Political manifestos are election platforms, and political
scientists have spent decades turning them into numbers. The
numbers feed downstream models of voting, government formation,
coalition bargaining, and policy outcomes, and the established
pipeline is human reading: a trained coder or country expert reads
the platform and places the party on a handful of issue scales.
The supply of expert reads is finite, the documents are long, and
any automated stand-in has to compress before it can score. The
question we set up here is whether a tree of compressed states can
carry the same evidence the expert reading does, on the same
target.

Take the Benoit benchmark's economic-policy scale as a running
example. A trusted coder reads the full party platform and places it
on the seven-point scale; that placement is the oracle \(f^\ast\)
we want a learned tree to recover. To be useful for the same
target, a C-Tree state for a span has to carry whatever evidence
could move that placement: tax promises, spending commitments,
privatization claims, labor-market policy, welfare language, fiscal
constraints, and the qualifications about who benefits or pays.

A scalar score per leaf is too small a state for this work. Suppose
one section of a platform promises a tax cut for small businesses,
and a later section expands health spending. Suppose a third section
makes both conditional on private delivery and on a deficit rule.
The economic position of the platform turns on how those three
pieces fit together. A parent that only inherits two scalar child
scores cannot recover the qualification. A parent that inherits two
policy-evidence states can carry the tax position, the spending
position, and the condition that ties them together.

A C-Tree state is the working memory the tree keeps about a span.
The query that comes out at the end is a number, but the state has
to carry more than a number. In economic policy, a state may record
that a party supports tax cuts for small businesses, protects health
spending, and funds the package through privatization. In
environmental policy, it may carry both emissions targets and
exceptions for energy security. In decentralization, it may need to
distinguish administrative devolution from fiscal autonomy, regional
veto power, and country-specific institutional language. The score
is read off the root, so the state has to survive every merge before
that read is meaningful.

The three local laws name the failure modes. C1 (sufficiency) fails
when a leaf summary drops the evidence the rubric needs. C2
(idempotence) fails when a state drifts after a re-summarization
pass, for example when a conditional tax cut becomes unconditional.
C3 (merge consistency) fails when two child states are each fine on
their own and lose their interaction once merged. C3 is the
characteristic manifesto problem: meaning often sits in the relation
between sections, not in either section alone.

Decentralization is the stress dimension in the current results
because the evidence is less consistent across parties and countries
than economic-policy evidence is. A party can endorse local
administration while opposing fiscal devolution. It can support
regional identity while keeping central budget authority. It can
discuss federal reform in country-specific institutional terms. A
universal summary trained to serve all six dimensions can preserve
the shared evidence well and still blur these
decentralization-specific distinctions. The local diagnostic gap
for decentralization is exactly that signal.

This picture also explains why node labels matter. A full-document
expert score calibrates the root. A leaf or merge label checks
whether a smaller state still carries the same evidence for the same
dimension. Manifesto quasi-sentence codings are a natural fit
because they already mark local policy statements. They certify a
tree only when sampled or aggregated under a rule that targets the
same policy dimension as the root score.

Sections~\ref{sec:v8-ctree-math}--\ref{sec:v8-audit-certificate}
formalize this picture. Section~\ref{sec:v8-ctree-math} defines the
C-Tree and the local laws,
Section~\ref{sec:v8-objective} writes the training objective, and
Section~\ref{sec:v8-audit-certificate} develops the audit.
